% ============================================================
%  preamble.tex — Computational Optimal Transport 解説ノート共通設定
% ============================================================

% --- 基本パッケージ ---
\usepackage{amsmath,amssymb,amsthm,mathtools}
\usepackage[dvipdfmx]{hyperref}
\usepackage{pxjahyper}

% --- 図・表 ---
\usepackage{tikz}
\usetikzlibrary{decorations.pathreplacing,calc,arrows.meta,patterns,positioning}
\usepackage{pgfplots}
\pgfplotsset{compat=1.18}

% --- 定理環境 (tcolorbox) ---
\usepackage{tcolorbox}
\tcbuselibrary{theorems,breakable}

\newtcbtheorem[number within=section]{definition}{定義}{
  colback=blue!5, colframe=blue!40!black, fonttitle=\bfseries, breakable
}{def}
\newtcbtheorem[use counter from=definition]{proposition}{命題}{
  colback=green!5, colframe=green!35!black, fonttitle=\bfseries, breakable
}{prop}
\newtcbtheorem[use counter from=definition]{theorem}{定理}{
  colback=red!5, colframe=red!40!black, fonttitle=\bfseries, breakable
}{thm}
\newtcbtheorem[use counter from=definition]{remark}{注意}{
  colback=yellow!5, colframe=yellow!40!black, fonttitle=\bfseries, breakable
}{rem}
\newtcbtheorem[use counter from=definition]{example}{例}{
  colback=gray!5, colframe=gray!50!black, fonttitle=\bfseries, breakable
}{ex}
\newtcbtheorem[use counter from=definition]{setting}{設定}{
  colback=purple!5, colframe=purple!40!black, fonttitle=\bfseries, breakable
}{set}
\tcolorboxenvironment{proof}{
  colback=white, colframe=black!20,
  leftrule=2pt, rightrule=0pt, toprule=0pt, bottomrule=0pt,
  breakable
}

% --- アルゴリズム環境 ---
\newcounter{algcounter}[section]
\newtcolorbox{algorithm}[1]{
  colback=white, colframe=black!70,
  fonttitle=\bfseries,
  title={Algorithm~\refstepcounter{algcounter}\thealgcounter\label{#1}},
  boxrule=0.8pt
}

% --- 便利マクロ ---
\newcommand{\R}{\mathbb{R}}
\newcommand{\N}{\mathbb{N}}
\newcommand{\E}{\mathbb{E}}
\DeclareMathOperator*{\argmin}{arg\,min}
\DeclareMathOperator*{\argmax}{arg\,max}
\DeclareMathOperator{\diag}{diag}
\DeclareMathOperator{\tr}{tr}
\DeclareMathOperator{\KL}{KL}
\newcommand{\abs}[1]{\lvert #1 \rvert}
\newcommand{\norm}[1]{\lVert #1 \rVert}
\newcommand{\inner}[2]{\langle #1,\, #2 \rangle}
\newcommand{\defeq}{\overset{\mathrm{def}}{=}}

% --- 参考文献 ---
\usepackage[numbers,sort&compress]{natbib}
